\documentclass[hidelinks]{article}
\usepackage[a4paper, total={7in, 10in}]{geometry}
\usepackage[dvipsnames]{xcolor}
\usepackage{amsmath, amssymb, amsthm}
\usepackage{tikz}
\usepackage{tkz-euclide}
\usepackage[unicode]{hyperref}
\usepackage[all]{hypcap}
\usepackage{fancyhdr}
\usepackage{listings}
\usepackage{subcaption}
\usepackage{enumitem}
\usepackage{graphicx}

\lstset{basicstyle=\ttfamily\small,columns=fullflexible,frame=single,breaklines=true,showstringspaces=false}

\title{\textbf{Multipole Expansion Test}}
\author{Diego Juarez}
\date{March 25th, 2026}

\pagestyle{fancy}
\fancyhf{}
\fancyhead[L]{Multipole Expansion Test}
\fancyhead[R]{Jackson-style Practice}
\fancyfoot[C]{\thepage}

\begin{document}
\hypersetup{bookmarksnumbered=true}
\pagecolor{black}
\color{white}
\maketitle

\begin{Large}
\tableofcontents
\end{Large}
\pagebreak

\section{Source Notebook}
This problem set follows the notebook
\url{../notebooks/Multipole expansion in electrostatics.ipynb}
and focuses on monopole, dipole, quadrupole, symmetry, and validity of far-field approximations.

\section{Expansion Framework}
\subsection{Problem 1: Coulomb Kernel Expansion}
Starting from
\[
\Phi(\mathbf{r})=\frac{1}{4\pi\varepsilon_0}\int \frac{\rho(\mathbf{r}')}{|\mathbf{r}-\mathbf{r}'|}\,d^3r',
\]
derive the idea of expanding the kernel for $r\gg r'$. State the small parameter controlling the approximation.

\subsection{Problem 2: Monopole Term}
Show that the leading far-field term is
\[
\Phi_{\text{mono}}(\mathbf{r})=\frac{1}{4\pi\varepsilon_0}\frac{Q}{r},
\]
where
\[
Q=\int \rho(\mathbf{r}')\,d^3r'.
\]
Explain why this term vanishes for a neutral distribution.

\subsection{Problem 3: Dipole Term}
Show that the next term has the form
\[
\Phi_{\text{dip}}(\mathbf{r})=\frac{1}{4\pi\varepsilon_0}\frac{\mathbf{p}\cdot \hat{\mathbf{r}}}{r^2},
\]
with dipole moment
\[
\mathbf{p}=\int \mathbf{r}'\rho(\mathbf{r}')\,d^3r'.
\]

\subsection{Problem 4: Field of an Ideal Dipole}
Starting from
\[
\Phi(\mathbf{r})=\frac{1}{4\pi\varepsilon_0}\frac{\mathbf{p}\cdot \mathbf{r}}{r^3},
\]
compute the electric field of an ideal dipole.

\section{Quadrupole and Symmetry}
\subsection{Problem 5: Quadrupole Structure}
Define the quadrupole tensor
\[
Q_{ij}=\int \rho(\mathbf{r}')\left(3x_i'x_j'-r'^2\delta_{ij}\right)d^3r'.
\]
Explain why it is symmetric and traceless in the standard convention.

\subsection{Problem 6: Two Equal and Opposite Charges}
For charges $+q$ at $(0,0,a)$ and $-q$ at $(0,0,-a)$:
\begin{enumerate}[label=\alph*)]
\item Find the total charge.
\item Find the dipole moment.
\item Determine which multipole term is the leading nonzero contribution.
\end{enumerate}

\subsection{Problem 7: Simple Quadrupole Arrangement}
Consider four point charges arranged so that the total charge and dipole moment vanish. Propose one explicit configuration and show that the leading far-field term is quadrupolar.

\section{Symmetry Arguments}
\subsection{Problem 8: Spherical Symmetry}
Show that a spherically symmetric charge distribution has only a monopole term in its exterior potential.

\subsection{Problem 9: Reflection and Inversion Symmetry}
Explain how reflection symmetry or inversion symmetry can force certain multipole moments to vanish. Give one concrete example for each case.

\subsection{Problem 10: Axial Symmetry}
Discuss how axial symmetry simplifies the angular dependence of the multipole expansion. Why are Legendre polynomials natural in this setting?

\section{Canonical Examples}
\subsection{Problem 11: Uniformly Charged Sphere}
Find the exterior multipole expansion of a uniformly charged sphere and compare it with the exact potential.

\subsection{Problem 12: Neutral Distribution with Nonzero Dipole Moment}
Construct a simple continuous or discrete neutral distribution with nonzero dipole moment. Write its leading far-field potential.

\subsection{Problem 13: Neutral Distribution with Vanishing Dipole Moment}
Construct a simple configuration with zero total charge and zero dipole moment but nonzero quadrupole moment.

\section{Validity and Error}
\subsection{Problem 14: Region of Validity}
Explain precisely when a multipole expansion is expected to converge well. What goes wrong when the observation point lies inside or too near the source distribution?

\subsection{Problem 15: Error Decay}
If the monopole and dipole terms are retained, estimate how the neglected terms scale with distance. Give the answer in powers of $1/r$.

\section{Challenge Problems}
\subsection{Problem 16: Translating the Origin}
Show that the dipole moment depends on the choice of origin unless the total charge vanishes. Then explain why, for a neutral system, the dipole moment is origin independent.

\subsection{Problem 17: Angular Pattern}
For an ideal dipole oriented along the $z$-axis, write the angular dependence of the potential and field. Explain why the dipole field has strong directional structure unlike the monopole field.

\end{document}
